
\subsubsection{Entropy-Based Uncertainty Quantification}

Token-level entropy — the mean per-token Shannon entropy over the greedy output distribution — is a low-cost UQ signal derived from a single forward pass. Kadavath et al. [2022] study self-calibration in large language models, including a P(True) signal, on models from Anthropic; their analysis targets proprietary models and does not systematically measure AUROC on open-source models or short factual QA benchmarks. Xiao and Wang [2022] conduct a broader comparison of UQ methods on natural language understanding tasks, establishing that the choice of UQ method matters for classification tasks; they do not study generative hallucination detection.

Semantic entropy \citep{Kuhnetal2023} extends token entropy by computing uncertainty over semantic equivalence classes rather than token sequences. N stochastic responses are grouped into clusters using bidirectional NLI entailment (microsoft/deberta-large-mnli in the lorenzkuhn/semantic_uncertainty implementation), and entropy is computed over the cluster size distribution. Kuhn et al. (2023) report AUROC ≈ 0.78 on TriviaQA and Natural Questions with several LLMs, a substantial improvement over token entropy baselines. Their evaluation does not report NLI aggregation rates, and their benchmarks use responses that are typically longer and more syntactically varied than HaluEval-QA short factual QA responses.

\subsubsection{Consistency-Based Hallucination Detection}

SelfCheckGPT \citep{Manakuletal2023} measures the consistency of N stochastic generations relative to a greedy response using BERTScore, NLI, or n-gram overlap. The design assumption is that consistent generation indicates factual confidence, while inconsistent generation signals hallucination. Manakul et al. (2023) validate on WikiBio biography generation, reporting AUC-PR ≈ 0.85 for the BERTScore variant with N=5 samples. The benchmark-specific nature of this result — biography generation from Wikipedia — differs from short factual QA in both response length and label construction methodology.

Xiong et al. [2023] study verbalized uncertainty: how well LLMs express their own confidence in natural language. This is a distinct paradigm from the post-hoc signals studied here.

\subsubsection{Hallucination Benchmarks}

HaluEval \citep{Lietal2023} is a large-scale hallucination evaluation benchmark covering QA, dialogue, and summarization. The QA subset provides binary labels: each question-answer pair is paired with a ChatGPT-generated plausible-but-incorrect hallucinated answer. Li et al. (2023) note that the QA subset is the most reliable of the three HaluEval tasks. A relevant characteristic of the label construction is that ChatGPT-generated hallucinations are designed to be confident-sounding and internally coherent, a property that may interact adversely with consistency-based UQ methods.

\subsubsection{The Cross-Method Comparison Gap}

The three methods have been validated on separate benchmarks: semantic entropy on TriviaQA/NQ \citep{Kuhnetal2023}, SelfCheckGPT on WikiBio \citep{Manakuletal2023}, and token entropy as a baseline across multiple settings. This siloed validation means that published results cannot be used to select among methods for a new deployment scenario. The present work is, to the authors' knowledge, the first controlled comparison of all three methods on HaluEval-QA with the same LLM under matched inference budgets.

